\ExerciseSection

\begin{enumerate}
\def\labelenumi{\arabic{enumi})}
\tightlist
\item
  设 \(f(z) = z^n + a_1z^{n-1} + \cdots + a_{n-1}z + a_n\) 是一个 \(n\) 次复系数多项式，并令 \(f(z) = \prod_{i=1}^{n}(z - z_i)\) 。令 \(r_0 = \sup_i |z_i|\) 。
\end{enumerate}

\begin{enumerate}
\def\labelenumi{\alph{enumi})}
\tightlist
\item
  证明：若实数 \(r > 0\) 满足
\end{enumerate}

\[
r ^ {n} \leqslant | a _ {1} | r ^ {n - 1} + | a _ {2} | r ^ {n - 2} + \dots + | a _ {n - 1} | r + | a _ {n} |,
\]

则 \(r_0 \leqslant r\) ，并由此推出

\[
r _ {0} \leqslant \sup \left(\mathrm{I}, \sum_ {k = 1} ^ {n} | a _ {k} |\right).
\]

\begin{enumerate}
\def\labelenumi{\alph{enumi})}
\setcounter{enumi}{1}
\item
  设 \((\lambda_i)_{1 \leqslant i \leqslant n}\) 是由 \(n\) 个 \(>0\) 的实数组成的有限序列，且满足 \(\sum_{i=1}^{n} \lambda_i^{-1} = 1\) 。证明 \(r_0 \leqslant \sup_k (\lambda_k |a_k|)^{1/k}\) {[}利用 \(a\){]} 。
\item
  由 \(a\) ) 推出：若诸系数 \(a_i\) 都不为零，则
\end{enumerate}

\[
r _ {0} \leqslant \sup \left(2 | a _ {1} |, 2 \left| \frac {a _ {2}}{a _ {1}} \right|, \dots , 2 \left| \frac {a _ {n - 1}}{a _ {n - 2}} \right|, \left| \frac {a _ {n}}{a _ {n - 1}} \right|\right).
\]

\begin{enumerate}
\def\labelenumi{\alph{enumi})}
\setcounter{enumi}{3}
\tightlist
\item
  由 a) 推出
\end{enumerate}

\[
r _ {0} \leqslant | a _ {1} - 1 | + | a _ {2} - a _ {1} | + \dots + | a _ {n - 1} - a _ {n} | + | a _ {n} |
\]

{[}考虑多项式 \((z - 1)f(z)]\) 。由此证明：若诸 \(a_i\) 都是实数且 \(>0\) ，则

\[
r _ {0} \leqslant \sup \left(a _ {1}, \frac {a _ {2}}{a _ {1}}, \dots , \frac {a _ {n - 1}}{a _ {n - 2}}, \frac {a _ {n}}{a _ {n - 1}}\right).
\]

\begin{enumerate}
\def\labelenumi{\arabic{enumi})}
\setcounter{enumi}{1}
\tightlist
\item
  代数数是在有理数域 \(\mathbf{Q}\) 上为代数的复数。证明由全体实代数数组成的序域 \(\mathbf{B}\) 是一个极大序域，但它对于由 \(\mathbf{R}\) 诱导的拓扑并不完备{[}这一拓扑与 \(\mathcal{C}_0(\mathbf{B})\) 相同；参看第四章，§ 3，习题 2 b){]}。
\end{enumerate}

¶ 3) a) 设 K 是一个极大序域，赋有拓扑 \(\mathfrak{G}_{0}(\mathbf{K})\) ，该拓扑与其域结构相容（第四章，§ 3，习题 2）。设

\[
f (\mathbf {X}) = a _ {0} \mathbf {X} ^ {n} + a _ {1} \mathbf {X} ^ {n - 1} + \dots + a _ {n}
\]

是 K{[}X{]} 中在 K 内有一个单根 \(\alpha\) 的多项式。证明：对 K 中每个元素 \(\varepsilon > 0\) ，存在 K 中的 \(\eta > 0\) ，使得每个多项式 \(g(\mathbf{X}) = b_{0}\mathbf{X}^{n} + b_{1}\mathbf{X}^{n-1} + \cdots + b_{n}\) ，只要对一切 i 有 \(|a_{i} - b_{i}| \leqslant \eta\) ，就恰有一个单根 \(\beta\) 满足 \(|\beta - \alpha| \leqslant \varepsilon\) 。（把 f 与 g 按 X - \(\alpha\) 的幂展开。）给出 \(\eta\) 关于 \(a_{i}, \alpha\) 与 \(\varepsilon\) 的一个上界。

\begin{enumerate}
\def\labelenumi{\alph{enumi})}
\setcounter{enumi}{1}
\item
  推出：若 K 是极大序域，则它的完备化 \(\hat{K}\) （关于 \(\mathcal{G}_{0}(K)\) ；这是一个带有自然序结构的域：第四章，§ 4，习题 2）是极大序域。
\item
  设 \(K_{0}\) 是一个序域，并设 \(S = K_{0}((X))\) 是 \(K_{0}\) 上一个未定元的形式幂级数域；把 S 中 \(\geqslant 0\) 的元素取为 o 以及最低次项系数 \textgreater0 的形式幂级数，以此给 S 一个线性序。证明赋有拓扑 \(\mathcal{G}_{0}(S)\) 的 S 是完备的，但 S 不是极大序域（证明多项式 \(Y^{p} - X\) 在 S{[}Y{]} 中对每个整数 p \textgreater{} 1 都不可约）。
\end{enumerate}

\(d\) ) 在 \(c\) ) 的例子中取 \(\mathbf{K}_0 = \mathbb{R}\) 。设 \(\mathbf{K}\) 是 \(\mathbf{S}\) 的一个极大序代数扩张；证明 \(\mathbf{K}\) 在拓扑 \(\mathfrak{C}_0(\mathbf{K})\) 下不完备。（把 \(\mathbf{S}\) 嵌入以良序指数的形式幂级数所成的域 \(\mathbf{E}\) ，其序按与 \(\mathbf{S}\) 相同的方式给定，再把 \(\mathbf{K}\) 嵌入 \(\Omega\) —— \(\mathbf{E}\) 的一个极大序扩张 —— 。注意 \(\mathbf{E}\) 是完备的，并给出 \(\mathbf{K}\) 中一个柯西序列的例子，它收敛于一个元素 \(f \in \mathbf{E}\) ，该元素在 \(\mathbf{S}\) 上不是代数的；如此选取 \(f\) ，使得存在无穷多个 \(\mathbf{E}\) 到 \(\Omega\) 的一个代数闭包中的 S-同构，并且 \(f\) 在这些同构下的像两两不同。{]}

\begin{enumerate}
\def\labelenumi{\arabic{enumi})}
\setcounter{enumi}{3}
\tightlist
\item
  证明：\(f\) 若是拓扑域 \(\mathbf{C}\) 到 \(\mathbf{C}\) 的一个子域上的连续同态（必然是单射），则它或是恒等自同构，或是自同构 \(z \to \overline{z}\) 。 {[}注意必须 \(f(x) = x\) 对一切 \(x \in \mathbf{Q}\) 成立，并由此推出 \(f(x) = x\) 对一切 \(x \in \mathbf{R}\) 成立。{]}证明存在无穷多个 \(\mathbf{C}\) 到 \(\mathbf{C}\) 的异于 \(\mathbf{C}\) 本身的子域上的不连续同构。
\end{enumerate}

¶ 5) a) 设 K 是四元数体 H 的一个非交换子体；证明 K 的中心 Z 是 H 的中心 R 的一个子域。（注意 K 的每个元素与由 Z ∪ R 生成的域 L 的每个元素交换；若 Z ⊕ R ，则域 L 将是 H 的一个极大子域，而 K 将包含于 L 中，与假设矛盾。）

\begin{enumerate}
\def\labelenumi{\alph{enumi})}
\setcounter{enumi}{1}
\tightlist
\item
  推出：\(f\) （不必连续）若是 \(\mathbf{H}\) 到 \(\mathbf{H}\) 的一个子体上的同构，则它是形如 \(x \to axa^{-1}\) 的一个 \(\mathbf{H}\) 的内自同构。 {[}\(f\) 在 \(\mathbf{R}\) 上的限制是 \(\mathbf{R}\) 到 \(\dot{\mathbf{R}}\) 的一个子域上的同构，这由 \(a\) ) 得出；利用第四章 § 3 的习题 3。{]}
\end{enumerate}

